% !TeX root = ../papers/nc-p3.tex
\section{三軸表示の全係数と正則性の境界}
\label{sec:cross-product-parameters}
$L=\operatorname{diag}(a_1,a_2,a_3)$とし、前節の二写像を
\[
 B_\epsilon((r,u),(s,v))
   =(rv+\epsilon(Lu)\times v,\;su+\epsilon D(u,Lv))
\]
に広げる。係数について
$\Delta=a_1^2a_2^2+a_2^2a_3^2+a_3^2a_1^2$、$p_i=a_j^2+a_k^2$（巡回順）と書く。
偶位置の原点を$[1:0:0:0]$と呼ぶ。

\begin{theorem}[三軸族の全次数判定]\label{thm:cross-product-parameters}
標準Koszul AS正則である必要十分条件は$\Delta\ne0$である。
三係数が非零な正則点では反対階数10、最小周期2で、ちょうど二係数が非零なら
両対称チャートの反対階数8の正則点である。三係数が非零な正則点のうち、
共通点スキームが有限である必要十分条件は$p_1p_2p_3\ne0$であり、
このときスキームは被約な二十点である。
\end{theorem}

\begin{proof}
係数を不定元にした行列で$B_+^{\mathsf t} B_-$の偶位置対称性が成り立つ。
反対行列を四つの文字次数で分けると、次数零の四行四列は行列式$-4\Delta$、
残る三つの二行四列は定数の階数2の行列の$a_i$倍である。
以下の三次恒等式とともに\path{computations/verify_cross_product_parameters.py}で
係数ごとに再現できる。

まず$a_1a_2a_3\ne0$とする。二周期表示で、位置$p$からの積を文字で表し、
$e=(-1)^p$と書くと六関係は
\begin{equation}\label{eq:axial-reduction}
 0i=-a_j jk+e a_k kj,\qquad i0=-a_k jk-e a_j kj .
\end{equation}
ここで$0$は一文字である。三次の次元は全トーラス上で20である。
実際、$000$以外は零番文字を消して三文字の語にできる。文字次数$i\ne0$にある七語には
$j0k,k0j$の二通りの計算から二関係を得る。それぞれ
\[
 a_i(jij-kik)+e(a_j jji-a_k ikk),\quad
 e a_i(kik-jij)+a_k kki-a_j ijj
\]
で、二つは独立である。次数零には六順列と$000$があり、$i0i$の両側の計算は
$-a_k jki-e a_j kji+a_j ijk+e a_k ikj$を与える。
異なる二つの$i$では台が異なるので少なくとも二式は独立である。これが全体の上界20を与える。
下界は巡回Gramテンソルの四つの独立な一文字縮約が両位置の三次交叉に入ることによる。
一文字切断の全階数は、二つの軸における評価で各四文字の一次独立性から分かる。

四平方の係数で$q_{0,p}=(0,a_1,a_2,a_3)$とし、$q_{i,p}$の
$0,i,j,k$成分を順に
\[
 (-a_i,0,e a_k(a_i^2+a_j^2),-e a_j(a_i^2+a_k^2))
\]
と取る。前節と同じ四指標で$q_{h,p}v=\sigma_h(v)q_{h,p+1}$となる。
上記コードでは有理関数体で文字次数ごとの三次関係行列にこれらの式を追加し、
全てで階数11を保つことを調べる。三次商の次元はトーラス全体で一定と示したので、
各ファイバーの関係像は部分束であり、上の多項式恒等式は全ファイバーへ延びる。
四係数行列の行列式は$4\Delta^2$である。

四平方を割ったときの次元の上界は、三係数が非零なら常に$1,4,6,4,1,0$である。
二文字以上の語から式\eqref{eq:axial-reduction}で零番文字を除く。
$0ii$と$ii0$の計算から三文字全ての異なる並べ替えは互いに非零定数倍である。
$j0k$の計算から$jij=kik$である。従って三次は四次元以下である。
四次で三文字全てが現れ、既に平方がない語は、相異なる三文字の連続部分を並べ替えて
残る文字と隣り合わせにすれば零になる。残る交互語$ijij$は
$i j i=k j k$または$j i j=k i k$で第一または第二の添字を取り替えられ、
六つとも等しい。五次では四次の代表の末尾を第五文字と同じにして零とできる。

一般ファイバーの元代数の各次数の次元は標準値であることを先に確認する。
$(1,1,1)$がKoszulであり、二次関係と三次交叉はその近傍で束である。
$n$を固定し次数$n-1$まで平坦な近傍で、
$W_3\otimes A_{n-3}\to R\otimes A_{n-2}\to V\otimes A_{n-1}$を見る。
二写像の合成は零で基点では中間項で完全だから、階数の下半連続性で階数は近傍で一定となる。
最後の余核が$A_n$であり次数別帰納ができる。従ってどの特殊点でも各次元は
有理関数体での標準値以上である。

$\Delta\ne0$なら上の四正規元は四平方を張るので、
先の商の上界と正規性から各始点で多項式環上の高々$1,4,6,4,1$個の生成元があれば足りる。
この自由加群のHilbert級数は$(1-t)^{-4}$以下で、先の下界と一致する。
従って全次数の次元が標準で、両側とも四正規元の多項式環上で自由である。

ここからKoszul性も従う理由を記す。二頂点で折り畳み、各$q_i$は二頂点での和を指す。
四階テンソルの一文字・二文字の隣接切断は全階数である。
$k,V,R,W_3,kw$を係数空間とする右Koszul列の最初の二微分は表示によって完全
（$V\otimes A$の項での完全性）であり、最後の二微分は左の最初の二微分の
$A$双対になる。実際、$w$の$(1,3),(2,2),(3,1)$切断による完全対合で
一文字を移す両包含写像が互いに転置となる。
左の始点単純から$A$への$\operatorname{Hom},\operatorname{Ext}^1$はともに零である。
前者は$q_0$倍が単射だからである。後者について左増大イデアルからの左加群写像$f$を取る。
可換な二つの正規元$q_0,q_1$について$q_0 f(q_1)=q_1 f(q_0)$であり、
左多項式環上の自由性と互いに素な変数であることから$f(q_i)=q_i b$ ($i=0,1$)となる。
任意の正次数元$z$で$q_0$は中心的なので$q_0f(z)=z f(q_0)=q_0zb$、
従って$f$は右$b$倍に延びる。頂点別には冪等元で切ればよい。
これで右Koszul列の最上二項の正ホモロジーが消える。残る中央のホモロジーは
既に得たHilbert級数とEuler等式から零である。線型解消と完全対合から標準AS性が従う。

一方$\Delta=0$なら、$v_i=1/a_i$で$(s,v)=(1,v)$を取ると
$v\cdot v=0$、$Lv=(1,1,1)$となる。上側の値$rv+(Lu)\times v$は
$v^\perp$に入る。これが零のとき$Lu\in v^\perp$である。
実際、非零等方$v$では$w\times v\in kv$と$w\cdot v=0$が同値である。
従って$\sum u_i=0$となり、下側の値$(\sum u_i)(1,1,1)$も零である。
この奇文字の評価階数は2以下であり、定理\ref{thm:single-opposite-evaluation-ranks}で排除される。

ちょうど二係数が非零なら、例えば$a_i=0$として、$T$を$i$座標だけを負にする行列と置く。
$B_-((r,u),(-s,Tv))=\operatorname{diag}(T,-I)B_+((r,u),(s,v))$で両対称表示になる。
$0i$は偶側の反対核の相異なる因子で、二つの評価像は上下の相補三平面である。
対応する奇因子は$0,i$で、四評価とも階数3である。
上の三次元単位積の左乗法には非零な二つの軸回転$[a_j e_j]_\times,[a_k e_k]_\times$が現れ、
これらは可換でないので積は非結合的である。定理\ref{thm:transverse-factor-amalgam}を適用する。
階数は先のブロックから8である。一係数以下しか非零でないと、零係数の軸文字で
片側の評価が階数2以下となる。これで全パラメータを尽くした。

最後に三係数が非零な正則点の零点を調べる。$r=0$では$v$が$Lu$に比例し、
規格化後に$s u_i+p_i u_j u_k=0$となる。全$p_i\ne0$なら三座標軸と
$u_i^2$が$p_i$に比例する四点を得る。$p_i=0$なら$u_i=s=0$の直線を含む。
$r=1$では原点以外で$Lu\cdot v=v\cdot v=0$となる。
$u_1u_2u_3\ne0$と仮定して$w_i=a_i v_i/u_i$と置くと下側の式から
$w_i=c(2u_i^2-\sum u_j^2)$となる（$c\ne0$、さもなくば$v=0$）。
従って$Lu\cdot v=0$から$2\sum u_i^4=(\sum u_i^2)^2$となる。
差を四一次式に分解すれば符号変更後に一座標が他の二つの和である。
代入により$(a_i v_i)^2$が三つとも同じ非零値となり、
$v\cdot v=0$は$\Delta=0$を要して矛盾する。
例えば$u_k=0$で残りが非零のとき、上側から$v_i=-a_j u_j v_k$、
$v_j=a_i u_i v_k$で、下側と合わせると$u_i^2=u_j^2=-1/p_k$を得る。
一座標だけ非零の解はない。よって全$p_i\ne0$なら前節と同様に
$1+12+7=20$点で各奇核は直線であり、完全交叉の長さから被約である。
階数10のときの最小周期は前節と同じである。
\end{proof}

\begin{proposition}[原点を指定した同型と像の次元]\label{prop:cross-product-marked-isom}
$a_1a_2a_3\Delta\ne0$で原点を指定した二表示が添字保持同型である必要十分条件は、
三つの平方の順序を忘れた射影組$[a_1^2:a_2^2:a_3^2]$が一致することである。
原点を忘れた標準AS同型類の像も二次元である。
\end{proposition}
\begin{proof}
偶側の反対像は対称二乗の全体なので、その二つの基底変換は射影的に一致する。
固定した偶原点での奇核も両方一意である。三つの点空間の基底変換を
$(r,u)\mapsto(Ar+h(u),P u)$、$(s,v)\mapsto(Bs+k(v),Qv)$、
$(t,z)\mapsto(Ct+\ell(z),Tz)$と書く。
二つの全射の対による隣接分解の一意性から、$B_+$の標的の変換は両軸で
$AQ,BP$のブロックとなり、$B_-$の標的ではその逆転置の共通スカラー倍となる。
従って$P^{\mathsf t}P$はスカラーである。
上側の残る乗法は$h(u)I+[L'Pu]_\times=A Q[Lu]_\times Q^{-1}$であり、
跡から$h=0$となる。$L$が可逆なので$Q$は全ての歪対称行列を共役で歪対称へ送る。
$Q^{\mathsf t}Q$は従って三つの座標回転と可換で、スカラーに限られる。
$P=pO_V,Q=qO_U$とすると、$L'$は$O_U L O_V^{-1}$の非零定数倍であり、
$L^{\mathsf t}L$の固有値の射影組の一致が必要である。

逆に共通の置換で座標を並べ直せば、$L'=h_0 L S$と書ける
（$h_0\ne0$、$S$は対角符号行列）。共通置換は外積に置換符号を掛け、
$D$を保つため、$r$の符号で補正できる。
対角符号と倍率については偶変換を$(r,u)\mapsto(h_0\det(S) r,u)$、
両奇変換を$(s,v)\mapsto(h_0s,Sv)$と取る。
二つの六次元値が共通のブロック変換
$\operatorname{diag}(h_0\det(S) S,h_0 I)$で移り、計量の相似比は$h_0^2$である。
これで十分性を得る。一般の点スキームは二十点なので、印を忘れる写像のファイバーは
高々二十通りの原点の選択で尽くされる。印付きの二次元像は次元を保つ。
\end{proof}

\begin{theorem}[階数10の全特殊値で印を忘れた同型]\label{thm:cross-product-unmarked-isom}
$a_1a_2a_3\Delta\ne0$なら、原点を指定しない添字保持同型の条件も
三つの平方の順序を忘れた射影組の一致である。
\end{theorem}
\begin{proof}
まず全$p_i\ne0$のとき、偶位置の配置だけから不変量を回収する。
四つの座標点を頂点とし、残る十六点は各座標面に四点ずつ載る。
面$h$上の四点の座標平方は一つの射影行であり、非零座標に全ての符号を付けた点で尽くされる。
第零列の符号を変える対角基底変更（二乗前は$\sqrt{-1}$倍）により、四行は
対称な零対角行列$H$となり、$H_{0i}=p_i$、$H_{ij}=1$ ($1\le i<j\le3$)である。

三点以上を通る直線の本数は、各頂点で6、残る各点で3である。実際、一つの面の四符号点は
どの三つも独立で、その六対の結線には座標のうち一つの頂点が載る。
二頂点を結ぶ辺にさらに点はない。同じ面にある二点の直線がその面の外の点を通ることもない。
そこで三つの相異なる面の符号点だけを調べればよい。
頂点の置換後に残る面を無限遠面に取っても、面の規格化は同じ形に戻せる。
実際、対角基底変更と各行の規格化で$H_{ij}$を$t_i t_j H_{ij}$に変え、
第零以外の三頂点間の値を全て1にできる。新たな三値$p'_i$は
四頂点の三通りの対分割の辺積$H_{0i}H_{jk}$の並べ替えに共通の定数を掛けたものである。
従って$2\sum_{i<j}p'_ip'_j-\sum_i(p'_i)^2\ne0$は保たれる
（元の値は$4\Delta$）。アフィン座標で面$i$の点は第$i$座標だけが零で、
残り二座標は絶対符号を除いて共通の$m_i$、$m_i^2=-1/p'_i$となる。
これら三点が共線なら、アフィン依存の三係数$k_i$は和が零で、
三座標の式から$k_i m_i$は符号を除いて同じ非零値となる。
従って$\pm\sqrt{p'_1}\pm\sqrt{p'_2}\pm\sqrt{p'_3}=0$となり、上の非零条件に反する。

この本数で四頂点は配置から決まり、配置の同型は単項行列である。
同じ頂点順序で二つの対称行列表示が行・列のスカラー倍で対応すると、
対称性と全辺の非零性から倍率は共通定数を除き$t_i t_j$の形となる。
従って三つの対分割の辺積は、順不同の射影組として固有である。
これは$[p_1:p_2:p_3]$であり、$2a_i^2=p_j+p_k-p_i$により所望の不変量を得る。

次に$p_3=0,p_1p_2\ne0$とする。台は一本の直線$r=u_3=0$と、
孤立した十点（原点、第三軸、第一・第二座標面に各四点）からなる。
孤立点だけの三点直線は、原点と第三軸で各四本、他の各点では二本である。
実際、各四符号点の属する面には原点と第三軸があり、二点の結線のうち
これら二頂点を通るものが該当する。面の異なる孤立点三つにこのほかの結線はない。
従って二頂点が特定でき、これらと残る各点で張る二つの面を唯一の直線と交わらせて
他の二頂点も順序を除き復元できる。二つの四点面の平方行について、
第零・第三列の交叉比$H_{10}H_{23}/(H_{13}H_{20})=p_1/p_2$は
面または二頂点の交換による逆数を除き固有である。これで$[0:p_1:p_2]$を得る。
二つの$p_i$が零なら上の方程式から台に二本の直線があり、平方の組は
順序と共通倍率を除き$(1,1,-1)$しかない。各場合の十分性は印付きの同型による。
\end{proof}

\begin{corollary}[階数8の座標境界の同型]
三軸族の全正則点で、印なし同型は三係数の平方の順序を忘れた射影組の一致と同値である。
\end{corollary}
\begin{proof}
階数の相違で二つの部分は交わらない。残る$a_1=0$、$a_2=b,a_3=c$の場合を
先の両対称座標で比較する。両反対核は$\langle 01,23\rangle$であり、
四直線の交叉の四特異点が座標頂点を指定する。従って同型の各側の行列は単項行列で、
置換は対分割$01|23$を保つ。両側の平方同士の係数行列は
\[
 M=\begin{pmatrix}0&-1&1&1\\-1&0&-b^2&-c^2\\
 -1&b^2&0&-b^2-c^2\\-1&c^2&-b^2-c^2&0\end{pmatrix}.
\]
$b^2+c^2\ne0$なら零配置から両置換は一致する。
三通りの対分割$ij|kl$で$M_{ij}M_{ji}M_{kl}M_{lk}$を取ると
$[(b^2+c^2)^2:c^4:b^4]$が得られ、最初の分割が指定される。
$u=b^2/(b^2+c^2),v=c^2/(b^2+c^2)$は$u+v=1$、
$u-v=u^2-v^2$であるため、射影比$[b^2:c^2]$が交換を除き決まる。
$b^2+c^2=0$は零成分の数で区別されて一軌道である。
印付き十分性の符号・置換による構成は零係数を含めても成立する。
\end{proof}
